% Options for packages loaded elsewhere
\PassOptionsToPackage{unicode}{hyperref}
\PassOptionsToPackage{hyphens}{url}
\PassOptionsToPackage{dvipsnames,svgnames,x11names}{xcolor}
\documentclass[
]{article}
\usepackage{xcolor}
\usepackage[a4paper, hmargin=2.2cm, vmargin=2.4cm, headheight=14.5pt, footskip=0.9cm, includehead, includefoot]{geometry}
\usepackage{amsmath,amssymb}
\setcounter{secnumdepth}{-\maxdimen} % remove section numbering
\usepackage{iftex}
\ifPDFTeX
  \usepackage[T1]{fontenc}
  \usepackage[utf8]{inputenc}
  \usepackage{textcomp} % provide euro and other symbols
\else % if luatex or xetex
  \usepackage{unicode-math} % this also loads fontspec
  \defaultfontfeatures{Scale=MatchLowercase}
  \defaultfontfeatures[\rmfamily]{Ligatures=TeX,Scale=1}
\fi
\usepackage{lmodern}
\ifPDFTeX\else
  % xetex/luatex font selection
\fi
% Use upquote if available, for straight quotes in verbatim environments
\IfFileExists{upquote.sty}{\usepackage{upquote}}{}
\IfFileExists{microtype.sty}{% use microtype if available
  \usepackage[]{microtype}
  \UseMicrotypeSet[protrusion]{basicmath} % disable protrusion for tt fonts
}{}
\makeatletter
\@ifundefined{KOMAClassName}{% if non-KOMA class
  \IfFileExists{parskip.sty}{%
    \usepackage{parskip}
  }{% else
    \setlength{\parindent}{0pt}
    \setlength{\parskip}{6pt plus 2pt minus 1pt}}
}{% if KOMA class
  \KOMAoptions{parskip=half}}
\makeatother
\usepackage{color}
\usepackage{fancyvrb}
\newcommand{\VerbBar}{|}
\newcommand{\VERB}{\Verb[commandchars=\\\{\}]}
\DefineVerbatimEnvironment{Highlighting}{Verbatim}{commandchars=\\\{\}}
% Add ',fontsize=\small' for more characters per line
\newenvironment{Shaded}{}{}
\newcommand{\AlertTok}[1]{\textcolor[rgb]{1.00,0.00,0.00}{\textbf{#1}}}
\newcommand{\AnnotationTok}[1]{\textcolor[rgb]{0.38,0.63,0.69}{\textbf{\textit{#1}}}}
\newcommand{\AttributeTok}[1]{\textcolor[rgb]{0.49,0.56,0.16}{#1}}
\newcommand{\BaseNTok}[1]{\textcolor[rgb]{0.25,0.63,0.44}{#1}}
\newcommand{\BuiltInTok}[1]{\textcolor[rgb]{0.00,0.50,0.00}{#1}}
\newcommand{\CharTok}[1]{\textcolor[rgb]{0.25,0.44,0.63}{#1}}
\newcommand{\CommentTok}[1]{\textcolor[rgb]{0.38,0.63,0.69}{\textit{#1}}}
\newcommand{\CommentVarTok}[1]{\textcolor[rgb]{0.38,0.63,0.69}{\textbf{\textit{#1}}}}
\newcommand{\ConstantTok}[1]{\textcolor[rgb]{0.53,0.00,0.00}{#1}}
\newcommand{\ControlFlowTok}[1]{\textcolor[rgb]{0.00,0.44,0.13}{\textbf{#1}}}
\newcommand{\DataTypeTok}[1]{\textcolor[rgb]{0.56,0.13,0.00}{#1}}
\newcommand{\DecValTok}[1]{\textcolor[rgb]{0.25,0.63,0.44}{#1}}
\newcommand{\DocumentationTok}[1]{\textcolor[rgb]{0.73,0.13,0.13}{\textit{#1}}}
\newcommand{\ErrorTok}[1]{\textcolor[rgb]{1.00,0.00,0.00}{\textbf{#1}}}
\newcommand{\ExtensionTok}[1]{#1}
\newcommand{\FloatTok}[1]{\textcolor[rgb]{0.25,0.63,0.44}{#1}}
\newcommand{\FunctionTok}[1]{\textcolor[rgb]{0.02,0.16,0.49}{#1}}
\newcommand{\ImportTok}[1]{\textcolor[rgb]{0.00,0.50,0.00}{\textbf{#1}}}
\newcommand{\InformationTok}[1]{\textcolor[rgb]{0.38,0.63,0.69}{\textbf{\textit{#1}}}}
\newcommand{\KeywordTok}[1]{\textcolor[rgb]{0.00,0.44,0.13}{\textbf{#1}}}
\newcommand{\NormalTok}[1]{#1}
\newcommand{\OperatorTok}[1]{\textcolor[rgb]{0.40,0.40,0.40}{#1}}
\newcommand{\OtherTok}[1]{\textcolor[rgb]{0.00,0.44,0.13}{#1}}
\newcommand{\PreprocessorTok}[1]{\textcolor[rgb]{0.74,0.48,0.00}{#1}}
\newcommand{\RegionMarkerTok}[1]{#1}
\newcommand{\SpecialCharTok}[1]{\textcolor[rgb]{0.25,0.44,0.63}{#1}}
\newcommand{\SpecialStringTok}[1]{\textcolor[rgb]{0.73,0.40,0.53}{#1}}
\newcommand{\StringTok}[1]{\textcolor[rgb]{0.25,0.44,0.63}{#1}}
\newcommand{\VariableTok}[1]{\textcolor[rgb]{0.10,0.09,0.49}{#1}}
\newcommand{\VerbatimStringTok}[1]{\textcolor[rgb]{0.25,0.44,0.63}{#1}}
\newcommand{\WarningTok}[1]{\textcolor[rgb]{0.38,0.63,0.69}{\textbf{\textit{#1}}}}
\usepackage{longtable,booktabs,array}
\newcounter{none} % for unnumbered tables
\usepackage{calc} % for calculating minipage widths
% Correct order of tables after \paragraph or \subparagraph
\usepackage{etoolbox}
\makeatletter
\patchcmd\longtable{\par}{\if@noskipsec\mbox{}\fi\par}{}{}
\makeatother
% Allow footnotes in longtable head/foot
\IfFileExists{footnotehyper.sty}{\usepackage{footnotehyper}}{\usepackage{footnote}}
\makesavenoteenv{longtable}
\setlength{\emergencystretch}{3em} % prevent overfull lines
\providecommand{\tightlist}{%
  \setlength{\itemsep}{0pt}\setlength{\parskip}{0pt}}
% ============================================================
% 土力学速查手册 —— 学术黑白 A4 版式（xelatex / pandoc）
% 主题：严谨、黑白、无彩色、无装饰
% ============================================================
\usepackage[fontset=windows]{ctex}            % 中文排版（宋体正文/黑体标题，自动伪粗体）
% 把常用符号划入 CJK 字体（SimSun），避免落到拉丁字体缺字：
% 带圈数字 ①-⑩、几何符号 ◆■□、箭头 →、希腊字母 αφ、数学符号 ≈
\xeCJKDeclareCharClass{CJK}{"2460->"24FF}
\xeCJKDeclareCharClass{CJK}{"25A0->"25FF}
\xeCJKDeclareCharClass{CJK}{"2190->"21FF}
\xeCJKDeclareCharClass{CJK}{"0370->"03FF}
\xeCJKDeclareCharClass{CJK}{"2248->"2248}% ≈（SimSun 有；− 等保留拉丁字体）
\usepackage{amsmath, amssymb}                 % 数学符号
\usepackage{booktabs}                         % 三线表
\usepackage{longtable}                        % 跨页表格
\usepackage{array}
\usepackage{fancyhdr}
\usepackage{enumitem}
\setlist{itemsep=2pt, topsep=4pt, leftmargin=1.6em}

% ---------- 页眉页脚：简洁学术 ----------
\pagestyle{fancy}
\fancyhf{}
\fancyhead[C]{\small\kaishu 土力学速查手册}
\fancyfoot[C]{\small\thepage}
\renewcommand{\headrulewidth}{0.4pt}
\renewcommand{\footrulewidth}{0pt}
\fancypagestyle{plain}{
  \fancyhf{}
  \fancyfoot[C]{\small\thepage}
  \renewcommand{\headrulewidth}{0pt}
}

% ---------- 表格：黑白三线风格 ----------
\renewcommand{\arraystretch}{1.35}
\setlength{\aboverulesep}{2pt}
\setlength{\belowrulesep}{2pt}

% ---------- 其他 ----------
\raggedbottom
\usepackage{bookmark}
\IfFileExists{xurl.sty}{\usepackage{xurl}}{} % add URL line breaks if available
\urlstyle{same}
% fallback for those not using the hyperref driver hyperxmp:
\makeatletter
\@ifundefined{xmpquote}{\newcommand{\xmpquote}[1]{#1}}{}
\makeatother
\hypersetup{
  colorlinks=true,
  linkcolor={black},
  filecolor={Maroon},
  citecolor={black},
  urlcolor={black},
  pdfcreator={LaTeX via pandoc}}

\author{}
\date{}

\begin{document}

\section{土力学快速查阅手册（期末考试 /
考研版）}\label{ux571fux529bux5b66ux5febux901fux67e5ux9605ux624bux518cux671fux672bux8003ux8bd5-ux8003ux7814ux7248}

面向大学生期末考试与考研复习的土力学系统性速查手册。\textbf{分章节建立文档，公式完整析出}，考前可用于系统过一遍公式，考场可快速定位某一公式。

\begin{center}\rule{0.5\linewidth}{0.5pt}\end{center}

\subsection{项目结构}\label{ux9879ux76eeux7ed3ux6784}

{\def\LTcaptype{none} % do not increment counter
\begin{longtable}[]{@{}
  >{\raggedright\arraybackslash}p{(\linewidth - 2\tabcolsep) * \real{0.5000}}
  >{\raggedright\arraybackslash}p{(\linewidth - 2\tabcolsep) * \real{0.5000}}@{}}
\toprule\noalign{}
\begin{minipage}[b]{\linewidth}\raggedright
路径
\end{minipage} & \begin{minipage}[b]{\linewidth}\raggedright
说明
\end{minipage} \\
\midrule\noalign{}
\endhead
\bottomrule\noalign{}
\endlastfoot
\texttt{README.md}（根目录） & 项目总览 README（项目首页） \\
\texttt{markdown版本/} & 全部 Markdown 源文档（README、各章、总册） \\
\texttt{markdown版本/README.md} & 本说明文件（手册使用说明） \\
\texttt{markdown版本/00-全书公式速查总表.md} \textasciitilde{}
\texttt{08-地基承载力.md} & 各章 Markdown 源文档 \\
\texttt{markdown版本/土力学速查手册-总册.md} & 合并总册（封面 + 目录 +
各章全文，不含 README） \\
\texttt{PDF导出/} & 各章与总册的 PDF（学术黑白 A4，pandoc + xelatex
构建） \\
\texttt{LaTeX版本/} & 每个 Markdown 对应的独立可编译 \texttt{.tex}
源文件 \\
\texttt{latex\_log/} & LaTeX
编译过程文件（\texttt{.aux}/\texttt{.log}/\texttt{.toc}
等，自动生成，无需关注） \\
\texttt{.claude/pdf-build/} & 构建管线（版式模板、Lua
过滤器、构建脚本） \\
\end{longtable}
}

\subsubsection{LaTeX 版本说明}\label{latex-ux7248ux672cux8bf4ux660e}

\texttt{LaTeX版本/} 中每个 \texttt{.tex} 均由对应 Markdown
生成，内嵌本手册版式模板（ctex
中文字体、三线表、页眉页脚），\textbf{独立自包含}------直接用
\texttt{xelatex} 编译即可得到与 \texttt{PDF导出/} 一致版式的 PDF：

\begin{Shaded}
\begin{Highlighting}[]
\CommentTok{\# 在项目根目录执行}
\ExtensionTok{xelatex} \StringTok{"LaTeX版本/土力学速查手册{-}总册.tex"}
\end{Highlighting}
\end{Shaded}

\begin{itemize}
\tightlist
\item
  各章 \texttt{.tex} 对应单章 PDF；\texttt{土力学速查手册-总册.tex}
  已含居中封面与带页码目录。
\item
  修改 Markdown 后（在项目根目录）重新生成 LaTeX
  版本：\texttt{bash\ .claude/pdf-build/build\_tex.sh}；一键导出全部
  PDF：\texttt{bash\ .claude/pdf-build/build.sh}。
\end{itemize}

\begin{center}\rule{0.5\linewidth}{0.5pt}\end{center}

\subsection{使用说明}\label{ux4f7fux7528ux8bf4ux660e}

\subsubsection{目录导航}\label{ux76eeux5f55ux5bfcux822a}

{\def\LTcaptype{none} % do not increment counter
\begin{longtable}[]{@{}
  >{\raggedright\arraybackslash}p{(\linewidth - 4\tabcolsep) * \real{0.3333}}
  >{\raggedright\arraybackslash}p{(\linewidth - 4\tabcolsep) * \real{0.3333}}
  >{\raggedright\arraybackslash}p{(\linewidth - 4\tabcolsep) * \real{0.3333}}@{}}
\toprule\noalign{}
\begin{minipage}[b]{\linewidth}\raggedright
章节
\end{minipage} & \begin{minipage}[b]{\linewidth}\raggedright
文件
\end{minipage} & \begin{minipage}[b]{\linewidth}\raggedright
内容
\end{minipage} \\
\midrule\noalign{}
\endhead
\bottomrule\noalign{}
\endlastfoot
速查总表 & 00-全书公式速查总表 &
全部核心公式按主题一表汇总，考前最后冲刺首选 \\
第一章 & 01-土的物理性质与分类 &
三相比例指标及其换算、塑性指数、土的分类 \\
第二章 & 02-土的渗透性与渗流 & 达西定律、渗透系数、渗透力、流土与管涌 \\
第三章 & 03-土中应力计算 & 自重应力、基底压力、附加应力（角点法） \\
第四章 & 04-土的压缩性与地基沉降计算 & e-p
曲线、分层总和法、规范法、固结理论 \\
第五章 & 05-土的抗剪强度 & 库仑公式、莫尔-库仑准则、三轴/直剪试验 \\
第六章 & 06-土压力与挡土墙 & 静止土压力、朗肯土压力、库仑土压力 \\
第七章 & 07-土坡稳定分析 & 无黏性土坡、圆弧滑动法、瑞典条分法、毕肖普 \\
第八章 & 08-地基承载力 &
临塑荷载、临界荷载、太沙基极限承载力、深宽修正 \\
\end{longtable}
}

\subsubsection{快速入手路径}\label{ux5febux901fux5165ux624bux8defux5f84}

\begin{itemize}
\tightlist
\item
  \textbf{考前 2 小时}：直接看「00-全书公式速查总表」。
\item
  \textbf{逐章复习}：按目录顺序，每章以「本章导览 →
  核心公式（参数/单位/适用说明）→ 公式索引小结」结构展开。
\item
  \textbf{某个公式记不清}：用编辑器全文搜索关键词（如''固结度''\,``朗肯''），或在章末「本章公式索引」定位。
\end{itemize}

\subsubsection{公式渲染说明}\label{ux516cux5f0fux6e32ux67d3ux8bf4ux660e}

本手册公式使用 \textbf{LaTeX 数学语法}，在以下环境可正常显示：

\begin{itemize}
\tightlist
\item
  \textbf{Typora}：默认支持，打开即可见。
\item
  \textbf{Obsidian}：需开启「设置 → 编辑器 → 数学公式」。
\item
  ** VS Code + Markdown Preview Enhanced**：默认支持。
\item
  \textbf{GitHub} 网页预览：默认支持 KaTeX 渲染。
\end{itemize}

排版约定：行内公式用 \texttt{\$...\$}，独立公式用
\texttt{\$\$...\$\$}，关键结论用 \texttt{\textbackslash{}boxed\{\}}
高亮，多步推导用
\texttt{\$\$\textbackslash{}begin\{aligned\}...\textbackslash{}end\{aligned\}\$\$}。

不想在编辑器中预览公式时，可直接用 \texttt{LaTeX版本/} 中的
\texttt{.tex} 文件编译成 PDF（见「项目结构 · LaTeX 版本说明」）。

\subsubsection{编写维护约定（后续扩充规范）}\label{ux7f16ux5199ux7ef4ux62a4ux7ea6ux5b9aux540eux7eedux6269ux5145ux89c4ux8303}

\begin{enumerate}
\def\labelenumi{\arabic{enumi}.}
\tightlist
\item
  \textbf{文件命名}：\texttt{NN-章节主题.md}，NN 两位数字保持排序。
\item
  \textbf{公式编号}：不依赖自动编号，正文用「公式
  1-1、1-2\ldots」手工编号，首个数字为章号。
\item
  \textbf{公式三段式}：每个公式后必须有「参数说明」「单位」「适用条件/注意」三段，保证可独立阅读理解。
\item
  \textbf{新章节模板}：沿用「本章导览 → 核心公式 →
  本章公式索引小结」结构，并在「00-全书公式速查总表」中登记新公式。
\end{enumerate}

\begin{center}\rule{0.5\linewidth}{0.5pt}\end{center}

\emph{符号约定全局统一：见 00-全书公式速查总表
顶部的「〇、符号约定速查」（含 \(C_c\)、\(C_u\) 双义警示）。重力加速度取
\(g=10\,\mathrm{m/s^2}\)，水重度
\(\gamma_w=10\,\mathrm{kN/m^3}\)，多数教材与考试用值。}

\end{document}
